\section{Discussion}
\label{sec:discussion}

\subsection{Shape and circuit, regime by regime}

Table~\ref{tab:verdict} and Fig.~\ref{fig:synthesis} collect the results.
The Tesla valve fails in graphene in every regime, and for one reason. Its
asymmetry is an inertial loss: it needs a fluid, and a Reynolds number of
order ten. The ballistic electron gas is not a fluid, so the loops only add
resistance. The viscous electron liquid is one, but in graphene its
Reynolds number is capped by the Fermi velocity at values where the valve
gives a per cent or less, and the gate adds an asymmetry that is the same
for a straight strip. At terahertz frequencies the valve is long compared
with the plasmon wavelength and its resistance is two orders of magnitude
from an antenna. The one electronic Tesla valve with a large
diodicity~\cite{sarypov2026electron} works in another material, above a
threshold current, in a regime that is not the steady laminar flow of this
paper.

The funnel rectifies in every regime, and what does the rectifying differs.
In the electron fluid it is the circuit. In steady flow the shape leaves an
inertial asymmetry of less than $10^{-3}$ inside the drift limit at the
parameters of Table~\ref{tab:params}, the same with and without the gate,
while the connection of one contact to the gate
produces a per cent at $43$~mV in any channel. At terahertz frequencies the
gated funnel's rectification is that of its terminations: tying the drain
to the gate and driving the source produces $0.1$ to $0.3\AW$ in a funnel,
as it would in a plain channel of the same impedance, while the symmetric
drive that leaves only the shape produces a comparable response only near a
taper resonance at low temperature, and that response saturates first at
large signal.

In the ballistic gas it is the shape, together with the potential that the
bias lays over it, and the first condition for it to work at a frequency
$f$ is one of size:
\begin{equation}
L\lesssim0.4\,\frac{\vF}{f} .
\label{eq:criterion}
\end{equation}
This is the half-response point of Fig.~\ref{fig:length}(b), for a field
spread over the device; with the field computed the $0.55\um$ funnel is at
half response at $0.5\,\vF/L$. At $1\THz$ the criterion gives $0.4\um$.
The $6\um$ funnel on which the dc asymmetry of geometric diodes is best
documented is fifteen times too long, and its response at $1\THz$ exists
only if the potential drop is confined to the neck. A statement that
geometric diodes are or are not terahertz rectifiers has no content
without the length, and the dc asymmetry of a long device is no guide: the
best dc asymmetry at $\lambda=700\nm$ in Table~\ref{tab:short} belongs to
devices between $0.55$ and $2\um$, and only the shortest of them still
responds at $1\THz$ in a uniform field.

Length is necessary and not sufficient. The $0.55\um$ funnel meets the
criterion, and with an imposed potential it appears to give $0.7\AW$ per
absorbed watt whatever the profile. With the potential computed it gives
$0.12\AW$ at the room-temperature mean free path, because the bias does
not go where either imposed profile puts it: a quarter is spent at the
contacts and the rest is centred on the neck. This is a caution beyond the
present device. A uniform field is a common assumption of simulations of
geometric diodes; self-consistent electrostatics was already known to
roughly halve the static asymmetry~\cite{truccolo2024comprehensive}, and
here it reduces a static asymmetry of $1.56$ to $1.05$ and the terahertz
response sixfold. Three further conditions follow from
Sec.~\ref{sec:pic}. The mean free path is the least of them: from no
scattering to $100\nm$ the response per absorbed watt changes from $0.13$
to $0.10\AW$. The edges are the strictest: with half of the wall
reflections diffuse the response is lost. And the gate must be far: under
a gate at $30\nm$ a straight strip with its drain tied to the gate
rectifies twice the current of the funnel shape at $1\THz$, so that the
circuit takes over in the ballistic gas as it did in the fluid.
Shape-based rectification at terahertz frequencies, then, exists in one
of the three transport limits, in devices of half a micrometre and below
with specular edges and a distant gate, and at its best it equals the
circuit term of the plasmonic element at room temperature.

\begin{table*}
\caption{Rectification by geometry and regime. $A$: ballistic asymmetry at
$0.1$~V; $D$: asymmetry of the voltage in steady flow; $\RI$: current
responsivity into a $75\ \Omega$ source; $R_{\mathrm{abs}}$: per absorbed
watt, Eq.~\eqref{eq:rabs} for the imposed fields and
Eq.~\eqref{eq:computed} for the computed one. ``Shape only'' is the symmetric termination; DS
has the drain tied to the gate.}
\label{tab:verdict}
\begin{ruledtabular}
\begin{tabular}{L{0.15\textwidth}L{0.31\textwidth}L{0.24\textwidth}L{0.23\textwidth}}
Regime & Funnel & Classical Tesla valve & What limits it \\
\hline
Ballistic, dc & $A=1.6$ to $1.9$ clean; $1.1$ to $1.6$ at $\lambda=700\nm$, by length and field profile & $A=0.89$ to $1.08$, no systematic sign & reciprocity; needs the bias landscape \\
Viscous, dc, no gate & $|D-1|\leq10^{-2}$ within the drift limit & $10^{-3}$ to $3\times10^{-2}$ within the drift limit & drift speed below $0.3\,\vF$ caps $\Rey$ \\
Viscous, dc, gated & contacts exchanged: $<10^{-3}$; polarity reversed: $0.48\,V/U_{0}$ & the same polarity term; exchange $<10^{-3}$ & the polarity term is the gate, in any channel \\
Plasmonic, shape only & $10^{-4}$ to $0.13\AW$; $0.22\AW$ at taper resonance & $10^{-5}$ to $10^{-3}\AW$; $0.02\AW$ at resonance & plasmon damping; saturates first \\
Plasmonic, DS & $0.07$ to $0.34\AW$; $\eta<1\%$ & $0.02$ to $0.1\AW$ at $6$ to $11$~k$\Omega$ & $r^{2}$; gate overdrive; contacts \\
Ballistic, $1\THz$, $6\um$ & $\RI=0.05\AW$ (uniform field); $0.2$ to $0.45\AW$ (drop at neck) & none ($A=1$ at dc) & length: $fL/\vF=6$ \\
Ballistic, $1\THz$, $0.55\um$, imposed field & $R_{\mathrm{abs}}=0.7\AW$, either profile; $0.14$ to $0.42\AW$ with rough edges or $\lambda=100\nm$ & none & the field is not computed \\
Ballistic, $1\THz$, $0.55\um$, computed field & $R_{\mathrm{abs}}=0.12\AW$; $0.10\AW$ at $\lambda=100\nm$; $|R_{\mathrm{abs}}|\leq0.04\AW$ with half-diffuse edges; $\RI=0.012$ to $0.026\AW$ & not computed & contact drops; edges; under a $30\nm$ gate the circuit term is larger; $\eta\sim10^{-4}$ \\
\end{tabular}
\end{ruledtabular}
\end{table*}

\begin{figure*}
\includegraphics[width=\textwidth]{\figpath/fig8_synthesis.pdf}
\caption{Synthesis. (a) Current responsivity against frequency for every
mechanism: the plasmonic $1\um$ funnel with the Dyakonov-Shur termination
and with the symmetric drive that isolates the shape, the Tesla valve, the
$6\um$ ballistic funnel with a uniform field and with the potential drop
confined to the neck, all into a $75\ \Omega$ source, and the $0.55\um$
ballistic funnel with the field computed (gate at $300\nm$, shape-only
drive), for one device on the same source and per absorbed watt. Solid
lines: $100$~K parameters or the clean ballistic limit; dashed: $300$~K
or $\lambda=700\nm$. Points consistent with zero are not drawn. (b)
Noise-equivalent power of the simulated elements, which is Johnson
limited and assumes ideal coupling, against measured graphene terahertz
detectors (Table~\ref{tab:detectors}); the ballistic array is that of
Sec.~\ref{sec:pic}.}
\label{fig:synthesis}
\end{figure*}

\subsection{What the elements are good for}

Table~\ref{tab:detectors} and Fig.~\ref{fig:synthesis}(b) place the
simulated elements against measured graphene terahertz detectors. At room
temperature the best reported devices have noise-equivalent powers of $100$
to $130\pWrtHz$ referred to the optical power at
$0.4\THz$~\cite{generalov2017graphene,ludwig2024terahertz}, with an
estimate of $51\pWrtHz$ at $0.33\THz$ for an epitaxial bilayer
device~\cite{qin2017room}, and $80$ to
$160\pWrtHz$ referred to a diffraction-limited spot at $2.5$ to
$2.8\THz$~\cite{castilla2019fast,viti2020hbn}; most others lie between
$0.5$ and
$2$~nW/Hz$^{1/2}$~\cite{spirito2014high,zak2014antenna,bandurin2018dual,tong2015antenna},
the first devices were at tens of
nW/Hz$^{1/2}$~\cite{vicarelli2012graphene}, and four-terminal ballistic
rectifiers of encapsulated graphene have been reported at
$34\pWrtHz$ below $0.7\THz$~\cite{auton2017terahertz} and, at $3\THz$ and
$4$~K, at $20$ to $30\pWrtHz$~\cite{shi2026ballistic}; funnel rectifiers of
chemical vapor deposited graphene, with a mean free path below the device
size, are at $29$~nW/Hz$^{1/2}$~\cite{hemmetter2026enhanced}. The
Johnson-limited $60\pWrtHz$ of the DS-terminated $1\um$ funnel at $300$~K,
or $150\pWrtHz$ with $300\ \Omega\um$ contacts, is therefore where the best
measured devices already are, and its voltage responsivity of $50$~V/W is
within their range of $1$ to $100$~V/W. The idealised model holds no hidden
order of magnitude; the comparison supports the model rather than
promising an improvement. The matched array of $0.55\um$ ballistic funnels
is at $130\pWrtHz$ at $1\THz$ and room temperature, twice the figure of
the plasmonic element at that temperature and again where the best
measured room-temperature devices are; it assumes nineteen devices matched
to the antenna, specular edges and no parasitic capacitance, and it is
four times the value reported for the encapsulated ballistic
rectifiers~\cite{auton2017terahertz}. At low temperature, encapsulated devices reach
$0.2\pWrtHz$ at $10$~K and
$0.13\THz$~\cite{bandurin2018resonant,gayduchenko2021tunnel} and
$18\pWrtHz$ at $77$~K and $2.8\THz$~\cite{viti2020hbn}; the simulated $17$
to $27\pWrtHz$ at $100$~K, and $1\pWrtHz$ at low density, lie in the same
range. Three caveats apply. Nearly all reported values, like ours, are
computed from the Johnson noise of the measured resistance and not from
measured noise. The reference power differs between reports (total beam,
power delivered to the antenna, or power scaled to a diffraction-limited
spot), and values normalised to the active device area are not comparable
and are left out. And part of the measured response of graphene detectors
is photothermoelectric~\cite{cai2014sensitive,castilla2019fast,ludwig2024terahertz},
which the model does not contain. The room-temperature incumbent is better
than any of these: commercial zero-bias Schottky detectors have typical
noise-equivalent powers of $5$ to $27\pWrtHz$ from $0.22$ to
$1.1\THz$~\cite{vdi2026detectors}, which the simulated plasmonic funnel
reaches only at $100$~K and the simulated ballistic array does not reach.

\begin{table*}
\caption{Measured graphene terahertz detectors and the simulated element.
Responsivity and NEP as stated by the authors; the last column gives the
power to which they are referred: the power in a diffraction-limited spot,
the full optical beam power, the power delivered to the antenna, an
extrinsic figure, the waveguide input, or, for this work, the available
power with ideal coupling or the power absorbed by a matched array. NEPs are estimated from Johnson noise except the one marked with
a double dagger, which was measured; the commercial values are typical
specifications. Entries marked with a dagger were read from the abstract
only.}
\label{tab:detectors}
\begin{ruledtabular}
\begin{tabular}{lL{0.21\textwidth}cccL{0.14\textwidth}l}
Reference & Device and mechanism & $f$ (THz) & $T$ (K) & Responsivity & NEP (pW/Hz$^{1/2}$) & Power \\
\hline
\cite{vicarelli2012graphene} & FET, overdamped plasma wave & 0.3 & 300 & 0.10 to 0.15 V/W & $3\times10^{4}$ & spot \\
\cite{spirito2014high} & bilayer FET & 0.29 to 0.36 & 300 & 1.2 V/W & $2\times10^{3}$ & spot \\
\cite{zak2014antenna}$^{\dagger}$ & CVD FET, bow-tie antenna & 0.6 & 300 & 14 V/W & 515 & optical \\
\cite{generalov2017graphene} & CVD FET, bow-tie antenna & 0.4 & 300 & 74 V/W & 130 & optical \\
\cite{ludwig2024terahertz} & CVD FET, bow-tie antenna & 0.4 & 300 & 59 to 63 V/W & 101 to 114 & optical \\
\cite{bandurin2018dual} & encapsulated FET & 0.13 & 300 & 20 V/W & 600 & antenna \\
\cite{bandurin2018resonant} & encapsulated bilayer FET & 0.13 & 10 & $>3$ kV/W & 0.2 & antenna \\
\cite{tong2015antenna} & thermoelectric, patch antenna & 1.9 & 300 & 4.9 V/W & $1.7\times10^{3}$\,$^{\ddagger}$ & optical \\
\cite{castilla2019fast} & photothermoelectric p-n junction & 2.5 & 300 & 105 V/W & 80 & spot \\
\cite{viti2020hbn} & photothermoelectric and plasma wave & 2.8 & 300; 77 & 49; 180 V/W & 160; 18 & spot \\
\cite{auton2017terahertz}$^{\dagger}$ & four-terminal ballistic rectifier & up to 0.69 & 300 & 764 V/W & 34 & extrinsic \\
\cite{qin2017room}$^{\dagger}$ & epitaxial bilayer FET & 0.33 & 300 & 30 V/W & 51 & not verified \\
\cite{shi2026ballistic} & four-terminal ballistic rectifier & 3.0 & 4; 300 & 77; 6 V/W & 20 to 30 (4 K) & spot \\
\cite{hemmetter2026enhanced} & CVD four-terminal funnel rectifier & 0.49 to 0.68 & 300 & 0.77 V/W & $2.9\times10^{4}$ & optical \\
\cite{vdi2026detectors} & commercial zero-bias Schottky & 0.22 to 1.1 & 300 & 750 to 2200 V/W & 5 to 27 & waveguide \\
\hline
This work & $1\um$ funnel, DS, simulated & 0.3 to 1 & 300 & 0.09 A/W, 52 V/W & 60; 150 with contacts & available \\
This work & same & 0.3 to 1 & 100 & 0.21 to 0.34 A/W, 36 to 58 V/W & 17 to 27 & available \\
This work & same, $5\times10^{10}\percm$ & 0.3 & 100 & & 0.9 & available \\
This work & $0.55\um$ ballistic funnel, computed field, matched array, simulated & 1 & 300 & 0.12 A/W & 130 & absorbed \\
\end{tabular}
\end{ruledtabular}
\end{table*}

For energy conversion the conclusion does not depend on such comparisons,
and its framework is not new. The efficiency of a field-effect detector is
quadratic in the amplitude up to the gate swing and has its maximum, of the
order of ten per cent, below threshold~\cite{kachorovskii2013performance};
silicon transistor rectennas work there and have reached $13.6\%$ at
$263\GHz$~\cite{khan2022cmos}. What the present calculation adds is the
graphene case. A gapless monolayer has no subthreshold regime, the voltage
scale of its gated fluid is the gate overdrive, and the efficiency at an
amplitude equal to the overdrive is the pure number $r^{2}\Rant/\Rdc$ of
Eq.~\eqref{eq:eta}. For a funnel on a real antenna that number is of order
$10^{-2}$, and no density or gate distance in the hydrodynamic window
raises it; a low density or a thin gate moves the power at which the bound
is reached from a milliwatt to a microwatt, which is useful for matching a
detector's dynamic range to its signal, and nothing more. Realistic
contacts of $100$ to $300\ \Omega\um$ cost a further factor two to five.
The ballistic funnel has a lower voltage scale, the Fermi energy, and a
lower efficiency still: $10^{-4}$ at $13\ \mu$W per device, the largest
amplitude computed. A conversion efficiency of $58\%$ at
$0.9\THz$ has been predicted by simulation for a graphene
rectenna~\cite{dragoman2016graphene}, but its rectifier is a Schottky
junction between graphene and a semiconductor, not a geometric or plasmonic
element. Within the mechanisms studied here a graphene rectenna has no
harvesting regime.

A designer who wants a graphene terahertz detector from these shapes
should therefore read the results as follows. In a gated device of a
micrometre or more, use the funnel as a matched horn, not as a diode: its
value is the impedance it presents and the taper resonance. Terminate the
channel asymmetrically, since that is where the rectification comes from,
put the saturation knee at the expected power by choice of density and
gate distance, and spend the fabrication effort on the contacts, because a
resistive drain contact removes the asymmetry it is supposed to serve. If
the device is to rectify by its shape, make it short, by
Eq.~\eqref{eq:criterion}: half a micrometre for $1\THz$, leads included.
Spend the effort on the edges before the mobility, since a mean free path
of $100\nm$ costs a fifth of the response and half-diffuse edges cost all
of it; keep the gate farther away than the device is long, or accept that
the gate connection does the rectifying; and provide for the impedance,
since one such funnel is a kilohm. Expect no more than the circuit term
gives. Do not judge a geometric diode for terahertz use by its dc
asymmetry, nor by a simulation with an imposed field. Do not add Tesla
loops.

\subsection{Limitations}

There is no experiment here: the solvers are validated against analytic
limits and one measured number, and the agreement with that number depends
on the field profile and on the mean free path at the level of a factor of
several (Sec.~\ref{sec:dc}). Most idealisations favour the device. That is
the safe direction for the negative results, on the Tesla valve and on the
conversion efficiency, and the unsafe one for the positive result on the
short ballistic funnel, which computing the field reduced sixfold and
which remains an upper estimate.

The ballistic model of the $6\um$ device and of the length scaling
imposes its potential and has no displacement current. For the $6\um$
device the terahertz result depends entirely on the length of the drop,
and the scaling curve of Fig.~\ref{fig:length}(b) was confirmed with a
computed field for one device only, through its roll-off. For the
$0.55\um$ funnel the field is computed, in a model with limits of its
own. The dielectric is uniform. The contacts are ideal equipotential
strips that absorb and inject without resistance and without doping the
graphene beneath them; a quarter of the bias drops at the contact planes
within three grid cells, and since the rectified current goes roughly as
the square of what remains at the neck, a different contact would change
the result. The electrons are classical particles in a classical field,
with five to seven modes in the neck, and they interact through the
common field only; short-range electron-electron scattering enters
through the $100\nm$ bound. Two wall specularities were computed, and the
response is lost between them. The antenna is a voltage source, with its
$75\ \Omega$ entering through the computed admittance; the matched array
is an idealisation, and no array or antenna was modelled. Rectified
currents below about $0.05\ \mu$A are not resolved. Two results of
Sec.~\ref{sec:pic} are reported without an explanation: the reversal of
the strip's current near the plasma resonance at $3\THz$, where three
frequencies do not resolve the resonance, and the rise of the funnel's
current above the square law at $0.2$~V. No model of the crossover between
the ballistic and the fluid limit is offered: the two terahertz models
describe the same gate and nearly the same density, $1.4\times10^{12}$
and $10^{12}\percm$, in opposite limits of the electron-electron
scattering rate, and a real device at room temperature lies between
them.

The plasmonic and the gated steady-flow models use a local gate
capacitance, which is least accurate for the $175\nm$ neck under a $30\nm$
gate; they have ideal contacts except where stated, perfect coupling to an
antenna represented by its resistance, and Johnson noise only. They contain
no electron heating, and so none of the photothermoelectric response that
contributes to measured graphene detectors; the fluid description is
quantitative to about $0.3\THz$ at $100$~K and $1\THz$ at $300$~K, where
$\omega\tauee=1$, and indicative beyond. The low-bias shape term of the
gated steady flow is bounded, not resolved. Finally, two experiments on
electronic Tesla valves are outside what the models reproduce: the
rectification reported in graphene at low
current~\cite{geurs2020rectification} is not obtained from steady
hydrodynamics with or without the gate, and the large diodicity above
threshold in GaAs~\cite{sarypov2026electron} belongs to a regime of
unsteady, heated flow that was not computed.
